% !TEX root = ../../main.tex
\subsection*{Simulation protocol}

The rate-heterogeneous simulations use full alignments and paired formula evaluation. The five-taxon tree is a balanced four-taxon subtree \(T_A\) with all subtree branches set to 0.2 and a single taxon \(B\) attached by branch \(AB\). The 16-taxon tree is built from two balanced eight-taxon subtrees. Each alignment is analyzed under four scenarios. Two scenarios use the true topology: one with fixed true branch lengths and one with maximum-likelihood branch lengths. The other two use the maximum-likelihood inferred tree, without or with the Bonferroni threshold \( \alpha/(m_A m_B) \). Within each replicate, both formulas are evaluated on the same tree, branch length and target split.

The nucleotide simulations use the Tree-of-Life 16S rRNA GTR+F+G4 setting, implemented as the corresponding fixed four-category rate model used by the IQ-TREE integration \cite{minh2020iqtree}. The protein simulations use the Tree-of-Life LG+G4 setting. The growing-site grid uses \(100,250,500,1{,}000,2{,}000,4{,}000\) and \(10{,}000\) sites. It uses target branch lengths \(4,5,6,7.5,10\) and \(12\), with 100 replicates per condition. The 1,000-replicate grid uses \(100,250\) and \(1{,}000\) sites, target branch lengths \(4,5,6,7,8,9,10\) and \(12\). The branch-neighborhood analysis computes the same dominant and weighted statistics for the target branch and adjacent branches in the same simulated trees.
